\chapter{Cheilitis}

\section*{Definition}
Inflammation of the lips, presenting as dryness, cracking, erythema, and scaling of the vermilion border and surrounding skin. Chronic cases may involve fissuring at the oral commissures (angular cheilitis).

\section*{When to See a Doctor}
See your doctor if:
\begin{itemize}
  \item Lesions persist >2 weeks despite emollient use
  \item Severe pain, bleeding, or crusting interfering with eating
  \item Signs of secondary infection (purulent drainage, fever, spreading erythema)
  \item Suspicious non-healing ulcer or persistent scaly plaque (rule out squamous cell carcinoma)
\end{itemize}

\section*{How It Develops}
Irritation, allergy, sun exposure, dryness, or saliva can damage the lip barrier and cause inflammation, scaling, and fissuring. Angular cheilitis may involve saliva-associated maceration, a yeast or bacterial infection, denture factors, or another skin condition; infection is not present in every case.

\section*{Causes}
\begin{itemize}
  \item \textbf{Environmental:} Low humidity, wind, sun exposure (actinic cheilitis), cold weather
  \item \textbf{Irritant:} Lip-licking, mouth breathing, ill-fitting dentures, topical irritants
  \item \textbf{Infectious:} Candida albicans (angular cheilitis), herpes simplex virus, Staphylococcus aureus
  \item \textbf{Nutritional:} Riboflavin, iron, zinc, or B-vitamin deficiencies
  \item \textbf{Dermatologic:} Atopic dermatitis, contact dermatitis, psoriasis, lichen planus
  \item \textbf{Medication-related:} Isotretinoin, retinoids, chemotherapy agents
\end{itemize}

\section*{Related Symptoms}
\begin{itemize}
  \item Xerostomia (dry mouth)
  \item Glossitis (inflamed tongue)
  \item Perioral dermatitis
  \item Stomatitis
\end{itemize}
\textbf{Important:} Angular cheilitis in older adults often signals underlying denture-related stomatitis or nutritional deficiency and warrants full oral examination.

\section*{Medical Evaluation}
\begin{itemize}
  \item History: onset, triggers, medications, oral habits, dietary intake, denture use
  \item Inspection of lips, oral mucosa, and skin for concurrent dermatoses
  \item Swab or other testing when infection is suspected and the diagnosis is uncertain or treatment fails
  \item Patch testing if allergic contact dermatitis is considered
  \item Biopsy if actinic cheilitis or malignancy is suspected (non-healing ulcer, induration)
\end{itemize}

\section*{Treatment Options}
\begin{itemize}
  \item Emollients and barrier ointments (petrolatum, zinc oxide) as first-line for all types
  \item Barrier ointment is often sufficient for saliva-related angular cheilitis; an antifungal or other treatment is used when a clinician identifies the relevant cause
  \item Topical corticosteroid (hydrocortisone or desonide) for eczematous cheilitis
  \item Avoidance of identified allergens or irritants
  \item Nutritional supplementation if deficiency identified
  \item Cryotherapy, topical 5-fluorouracil, or photodynamic therapy for actinic cheilitis
\end{itemize}

\section*{Self-Care and Home Management}
\begin{itemize}
  \item Apply plain petrolatum or another bland, fragrance-free emollient frequently and before bed
  \item Avoid licking the lips and scented or flavored products; use lip protection with broad-spectrum sunscreen outdoors
  \item Use a humidifier in dry indoor environments
  \item For angular cheilitis that persists or recurs, seek assessment rather than repeatedly self-treating with antifungal or steroid creams
\end{itemize}

\section*{References}
\begin{thebibliography}{99}
\bibitem{cheilitis:ref1} Lugovic-Mihic L, Pilipovic K, Crnaric I, et al. ``Cheilitis: A Clinical Review.'' Acta Dermatovenerol Croat. 2020;28(2):93--104.
\bibitem{cheilitis:ref2} Scully C. ``Oral and Maxillofacial Medicine.'' 3rd ed. Churchill Livingstone, 2013.
\bibitem{cheilitis:ref3} Rosen T, Stone MS, Chuh A, et al. ``Cheilitis: A Comprehensive Review.'' J Am Acad Dermatol. 2022;86(4):757--768.
\bibitem{cheilitis:ref4} Nevill C, Campbell J, Day D. ``Actinic Cheilitis: A Systematic Review of Treatment Modalities.'' J Oral Maxillofac Surg. 2023;81(5):587--596.
\bibitem{cheilitis:ref5} Stoopler ET, Sollecito TP. ``Oral Mucosal Diseases.'' Med Clin North Am. 2014;98(6):1283--1303.
\end{thebibliography}
